


UniZ handles two types of campus exit permissions:

- **Outpass** — multi-day leave (overnight or longer). Requires full multi-level approval before you can leave.
- **Outing** — short-duration leave on the same day. Used for trips that start and end on the same day (e.g., going to town for a few hours).

<Note>
  You can only have one pending request at a time. You cannot submit a new
  outpass or outing while a previous request is still awaiting approval.
</Note>

---

\subsection{Request an outpass}

Use an outpass for any leave that involves an overnight stay outside campus.


\section{Instruction Step}

  
\section{Instruction Step}

    Go to your profile and select the **Permissions** tab. Tap **Request Outpass**.
  
  
\section{Instruction Step}

    Write a clear reason for your leave (e.g., "Emergency visit to home due to health issues"). This is reviewed by your caretaker, warden, and the SWO/Dean/Director in sequence.
  
  
\section{Instruction Step}

    Select your **From** date (when you plan to leave) and **To** date (when you plan to return). The return date must be on or after the departure date.
  
  
\section{Instruction Step}

    Tap **Request Outpass**. If successful, the system creates a new request at the `caretaker` approval level.

    \begin{lstlisting}[language=json]
    POST /requests/outpass
    {
      "reason": "Emergency visit to home due to health issues",
      "fromDay": "2026-02-10T09:00:00Z",
      "toDay": "2026-02-15T18:00:00Z"
    }
    
\end{lstlisting}

    **Response:**

    \begin{lstlisting}[language=json]
    {
      "success": true,
      "data": {
        "id": "req-789",
        "currentLevel": "caretaker",
        "isApproved": false,
        "createdAt": "2026-01-31T10:00:00Z"
      }
    }
    
\end{lstlisting}

  
</Steps>

---

\subsection{Request an outing}

Use an outing for same-day, short-duration trips that don't require an overnight stay.


\section{Instruction Step}

  
\section{Instruction Step}

    Go to your profile and select the **Permissions** tab. Tap **Request Outing**.
  
  
\section{Instruction Step}

    Write a brief reason (e.g., "Buying groceries").
  
  
\section{Instruction Step}

    Select your **From** time and **To** time. Both times apply to today's date.
  
  
\section{Instruction Step}

    Tap **Request Outing**.

    \begin{lstlisting}[language=json]
    POST /requests/outing
    {
      "reason": "Buying groceries",
      "fromTime": "2026-02-10T17:00:00Z",
      "toTime": "2026-02-10T20:00:00Z"
    }
    
\end{lstlisting}

  
</Steps>

---

\subsection{Approval workflow}

Both outpass and outing requests go through a multi-level approval chain before you can leave campus.

\begin{lstlisting}[language=]
You submit → Caretaker → Warden → SWO / Dean / Director (final approval)

\end{lstlisting}

| Level | Role                  | Action                                        |
| ----- | --------------------- | --------------------------------------------- |
| 1     | Caretaker             | Reviews and forwards to warden                |
| 2     | Warden                | Reviews academic schedule and forwards to SWO |
| 3     | SWO / Dean / Director | Grants final approval                         |

The `currentLevel` field in your request shows which level is currently reviewing it. Each approver may add a comment.

<Info>
  Caretakers are gender-specific — male caretakers review male students'
  requests, and female caretakers review female students' requests. The system
  routes your request automatically based on your profile.
</Info>

---

\subsection{Track your requests}

Open the **Permissions** tab on your profile to see your full request history. Each entry shows:

- Request type (outpass or outing)
- Date and time range
- Your reason
- Current status

\subsubsection{Request statuses}

| Status       | Meaning                                                |
| ------------ | ------------------------------------------------------ |
| **Pending**  | Still moving through the approval chain                |
| **Approved** | Fully approved — you may leave campus                  |
| **Rejected** | Declined by one of the approvers                       |
| **Expired**  | The request was not acted on before its departure date |

---

\subsection{After approval}

Once your request reaches `isApproved: true`:

- Security at the gate will see your approved request and allow you to exit.
- Your campus presence status (`is_in_campus`) is updated to `false` when you check out at the gate.
- When you return, security marks you back in and your status returns to `true`.

<Warning>
  Do not attempt to leave campus before your request is fully approved. Security
  verifies each student against the system — unapproved exits are recorded.
</Warning>

---

\subsection{If your request is rejected}

When an approver rejects your request, you'll see a rejection status and any comment they added. Common reasons include:

- Upcoming examinations or assessments
- Incomplete attendance records
- Insufficient reason provided
- Disciplinary holds

If you believe a rejection was in error, speak to your caretaker or warden directly.

After a rejection, the pending request clears and you can submit a new request if needed.


\section{Deep Technical Analysis}
The underlying system architecture for this module is built upon the \textbf{UniZ Microservices Framework}. 
This design ensures that each functional unit (Auth, Profile, Academics) remains isolated, preventing cascading failures across the university ecosystem.

\subsection{Operational Hardening}
Deployments are managed via K3s (Kubernetes) and containerized using high-performance Alpine Linux Docker images. 
Resource management is critical; for instance, the documentation service itself is provisioned with 2GB of RAM to ensure the responsive rendering of complex documentation graphs and state-management components.

\subsection{Enterprise Security Topology}
Every request entering the UniZ gateway is subjected to an aggressive JWT validation layer. 
Permissions are granular, mapping directly onto the RBAC (Role-Based Access Control) matrix defined in the core system specifications.

\subsection{Data Governance}
Prisma-driven data access ensures type safety and predictable query performance. 
We utilize PostgreSQL connection pooling to handle concurrent requests from the student portal and administrative dashboards simultaneously.
